% 05-isolation.tex
\chapter{Security and Isolation}
\label{ch:isolation}

\section{Isolation Framework}
\label{sec:isolation-framework}

Doki's isolation framework provides 12 levels of process isolation,
from hardware-enforced microVMs to lightweight syscall interception.
The framework is implemented as a registry of runtime modules, each
providing a different isolation mechanism. At startup, the daemon
probes the host system and selects the strongest available runtime.

\section{Isolation Levels}
\label{sec:isolation-levels}

The following table summarizes the 12 isolation levels, their
mechanisms, and their applicability to different environments.

\begin{table}[htbp]
  \centering
  \caption{Doki isolation levels with their mechanisms and
   applicable environments.}
  \label{tab:isolation-levels}
  \small
  \begin{tabularx}{\textwidth}{l l l X}
    \toprule
    \textbf{Level} & \textbf{Runtime} & \textbf{Mechanism} & \textbf{Applicable} \\
    \midrule
    1  & microVM    & Hardware virtualization (Firecracker/crosvm) & Cloud, server \\
    2  & gVisor     & User-space kernel (Sentry) & Server, restricted kernel \\
    3  & Namespaces & Kernel namespaces + cgroups & Linux with root \\
    4  & Chroot     & Filesystem isolation & Linux with root \\
    5  & Proot      & ptrace syscall interception & Android, rootless \\
    6  & PKdroid    & Android-specific proot & Android/Termux \\
    7  & QEMU-user  & User-mode emulation & Cross-architecture \\
    8  & FEX-Emu    & x86 emulation on ARM & x86 binaries on ARM \\
    9  & Legacy32   & 32-bit compatibility & 32-bit binaries \\
    10 & Sysbox     & Rootless container runtime & Server, rootless \\
    11 & WASM       & WebAssembly sandbox & Any platform \\
    12 & Native     & Direct execution & Trusted environments \\
    \bottomrule
  \end{tabularx}
\end{table}

\section{Seccomp Filtering}
\label{sec:seccomp}

Doki uses seccomp (secure computing mode) to restrict the system calls
available to container processes. The default seccomp profile allows
approximately 80 syscalls, including modern interfaces such as
\texttt{io\_uring\_*}, \texttt{pidfd\_*}, \texttt{rseq},
\texttt{userfaultfd}, \texttt{copy\_file\_range}, and
\texttt{landlock\_*}.

\subsection{Default Deny List}
\label{sec:default-deny}

The following syscalls are denied by default, regardless of the
isolation level:

\begin{codebox}[title={Syscalls denied by default seccomp profile}]{text}
init\_module, finit\_module, delete\_module,
kexec\_load, kexec\_file\_load,
iopl, ioperm, kcmp,
process\_vm\_readv, process\_vm\_writev
\end{codebox}

These syscalls are blocked because they can be used to load kernel
modules, modify I/O ports, or access the memory of other processes,
all of which could compromise the host system.

\subsection{Custom Profiles}
\label{sec:custom-profiles}

Users can customize the seccomp profile through the container
configuration:

\begin{codebox}[title={Custom seccomp configuration}]{json}
\{
  "SecurityOpt": \{
    "seccomp": \{
      "defaultAction": "SCMP\_ACT\_ERRNO",
      "architectures": ["SCMP\_ARCH\_X86\_64"],
      "syscalls": [
        \{
          "names": ["read", "write", "open", "close"],
          "action": "SCMP\_ACT\_ALLOW"
        \}
      ]
    \}
  \}
\}
\end{codebox}

\section{Linux Capabilities}
\label{sec:capabilities}

Doki manages Linux capabilities through the container configuration.
By default, containers run with a reduced capability set that excludes
dangerous operations such as \texttt{CAP\_SYS\_MODULE},
\texttt{CAP\_SYS\_RAWIO}, and \texttt{CAP\_SYS\_ADMIN}.

\section{AppArmor}
\label{sec:apparmor}

When available on the host system, Doki applies AppArmor profiles to
container processes. AppArmor provides mandatory access control that
restricts file access, network operations, and capability usage beyond
what seccomp alone provides.

\section{User Namespaces}
\label{sec:user-namespaces}

On systems that support user namespaces, Doki maps container user IDs
to unprivileged host user IDs. This provides an additional isolation
layer even when running as a non-root user, as container processes
cannot access files owned by the host user outside their mapped range.

\section{Isolation Level Selection}
\label{sec:level-selection}

The runtime registry in \texttt{pkg/runtime/registry.go} implements
the following selection algorithm:

\begin{enumerate}[leftmargin=*, itemsep=3pt]
  \item Probe for microVM support (check for Firecracker/crosvm
        binaries and KVM availability)
  \item Check for gVisor (\texttt{runsc}) availability
  \item Try unshare with all six namespace types
  \item Attempt chroot (requires root or \texttt{chroot} capability)
  \item Fall back to proot (always available on Linux)
  \item If on Android, try PKdroid-specific optimizations
  \item As last resort, use native execution
\end{enumerate}

The user can override automatic selection with the
\texttt{--runtime} flag:

\begin{codebox}[title={Forcing a specific runtime}]{bash}
doki run --runtime proot alpine:3.18 sh
doki run --runtime native ubuntu:22.04 bash
\end{codebox}
